\section*{Hoofdstuk \ref{ch:g6:producing-our-food} --- Ons voedsel produceren}

\begin{solution}{exo:g6:producing-our-food:1}
De \omterm{def:g6:producing-our-food:farming}{akkerbouw} beheert \omterm{def:g5:food-webs:roles}{producenten} --- \omterm{def:g1:plants-around-us:plant}{planten} die om hun \omterm{def:g6:plants-are-producers:matter}{organische
materie} geteeld worden. De \omterm{def:g6:producing-our-food:farming}{veehouderij} beheert \omterm{def:g5:food-webs:roles}{consumenten} --- \omterm{def:g1:animals-around-us:animal}{dieren} die
om hun productie gehouden worden: \omterm{def:g2:animals-grow-up:born}{melk}, vlees, eieren, wol.
\end{solution}

\begin{solution}{exo:g6:producing-our-food:2}
Graan: het \omterm{def:g2:from-seed-to-plant:seed}{zaad} van de tarweplant (haar ingepakte voorraad). Aardappel:
de \omterm{ex:g4:plants-without-seeds:bulbtuber}{knol} van de aardappelplant. Appel: de \omterm{prop:g3:flowers-fruits-seeds:transform}{vrucht} van de appelboom. \omterm{def:g2:animals-grow-up:born}{Melk}:
het product van de moederzorg van de koe. \omterm{def:g2:animals-grow-up:born}{Ei}: de ingepakte kraamkamer van
de kip.
\end{solution}

\begin{solution}{exo:g6:producing-our-food:3}
De omzetting van een grondstof door etende \omterm{def:g6:producing-our-food:fermentation}{microben}. Brood (grondstof:
bloemdeeg, werker: \omterm{ex:g6:producing-our-food:bread}{gist}) en yoghurt (grondstof: \omterm{def:g2:animals-grow-up:born}{melk}, werkers:
zuurmakende \omterm{def:g6:producing-our-food:fermentation}{microben}); kaas en azijn voldoen ook.
\end{solution}

\begin{solution}{exo:g6:producing-our-food:4}
De \omterm{ex:g6:producing-our-food:bread}{gist} --- een eencellige \omterm{def:g6:decomposers-and-soil:decomposers}{schimmel} --- voedt zich met de suikers van de
\omterm{def:g3:flowers-fruits-seeds:parts}{bloem} en geeft een gas af; het rekbare deeg vangt het gas als bellen en
blaast zich op. De hitte van de oven maakt een eind aan het werk van de
\omterm{ex:g6:producing-our-food:bread}{gist} en zet het gerezen kruim vast, met gaten en al.
\end{solution}

\begin{solution}{exo:g6:producing-our-food:5}
Het zuur dat de \omterm{def:g6:producing-our-food:fermentation}{microben} afgeven terwijl ze zich met de melksuiker
voeden: dat zuur dikt de \omterm{def:g2:animals-grow-up:born}{melk} in en laat haar stollen, en de frisse zure
smaak is het zuur zelf.
\end{solution}

\begin{solution}{exo:g6:producing-our-food:6}
Zouten, koelen (de koelkast), drogen --- en ook koken: elk daarvan
ontzegt ongewenste \omterm{def:g6:producing-our-food:fermentation}{microben} een werkomstandigheid (vocht, warmte),
precies zoals kou een composthoop vertraagt.
\end{solution}

\begin{solution}{exo:g6:producing-our-food:7}
Water: de boer vertrouwt op regen (en brengt die soms zelf). Licht:
zaaien op open \omterm{prop:g5:people-and-nature:managed}{akkers}, zo gespreid dat elke \omterm{def:g1:plants-around-us:plant}{plant} haar deel krijgt.
Warmte: het zaaien afgestemd op het warme seizoen. Lucht: open grond,
losgemaakte bodem. \omterm{prop:g4:what-plants-need:needs}{Minerale zouten}: mest en aanvulling na de oogst.
\end{solution}

\begin{solution}{exo:g6:producing-our-food:8}
(1) Organisme en stadium: koeienmelk --- een product van een \omterm{def:g5:food-webs:roles}{consument}.
(2) Plaats in de kringloop: product van een \omterm{def:g5:food-webs:roles}{consument}, door \omterm{def:g6:producing-our-food:fermentation}{microben}
afgemaakt. (3) Grondstof \omterm{def:g2:animals-grow-up:born}{melk}; de werkers maken haar zuur en laten haar
stollen tot wrongel, en rijpen daarna de geperste kaas maandenlang ---
gaten, korst en geur zijn hun logboek. (4) Spoor: \omterm{def:g2:animals-grow-up:born}{melk} van de koe, koe
van gras, gras van licht.
\end{solution}

\begin{solution}{exo:g6:producing-our-food:9}
\omterm{def:g6:producing-our-food:fermentation}{Microben} zijn levende werkers met \omterm{def:g6:exploring-our-environment:conditions}{omstandigheden}: kou legt hun eten
vrijwel stil. In de koelkast beschermt dat het voedsel (de bedervers
liggen stil); in de winterhoop vertraagt het de compost (de \omterm{def:g6:decomposers-and-soil:decomposers}{afbrekers}
liggen stil). Het beginsel: afbraak en bederf lopen op het tempo dat de
\omterm{def:g6:exploring-our-environment:conditions}{omstandigheden} toelaten.
\end{solution}

\begin{solution}{exo:g6:producing-our-food:10}
``Microbe'' is de naam voor een enorme menigte heel verschillende
piepkleine wezens. Sommige veroorzaken op de verkeerde plek ziekte --- daar richt
het handen wassen zich op; andere, gekozen en beheerd onder beheerste
\omterm{def:g6:exploring-our-environment:conditions}{omstandigheden}, laten brood rijzen en \omterm{def:g2:animals-grow-up:born}{melk} stollen. Vijand of werknemer
is niet de aard van de microbe maar de combinatie van \omterm{def:g5:biodiversity:species}{soort}, plaats en
beheer.
\end{solution}

\begin{solution}{exo:g6:producing-our-food:11}
De ham ging via een varken: zijn gewicht kostte vele malen zijn gewicht
aan voer, plus het verbrande deel van het varken, plus land en water ---
de tol van de consumentenverdieping. De tarwe van het brood kwam één
verdieping vanaf het licht. Het broodje weggooien verspilt allebei de
boekhoudingen; die van de ham is de langste.
\end{solution}

\begin{solution}{exo:g6:producing-our-food:12}
Ze beheerden de \emph{\omterm{def:g6:exploring-our-environment:conditions}{omstandigheden}} van de werkers: warmte voor de hoek
waar het deeg stond, het bewaarde en gevoede desem, zout en timing die
door generaties proberen afgesteld waren. De recepten werkten omdat de
juiste \omterm{def:g6:exploring-our-environment:conditions}{omstandigheden} aanhouden de juiste \omterm{def:g6:exploring-our-environment:components}{levende wezens} laat gedijen ---
of iemand nu weet dat die \omterm{def:g6:exploring-our-environment:components}{levende wezens} er zijn of niet. Ze hielden
blind \omterm{def:g6:producing-our-food:fermentation}{microben}, en goed.
\end{solution}

\begin{solution}{pb:g6:producing-our-food:1}
\textbf{1.} Elke korrel is een \omterm{def:g2:from-seed-to-plant:seed}{zaad}: een piepklein tarweplantje met een
\omterm{def:g3:animals-through-seasons:active}{voedselvoorraad} erbij, door de moederplant gemaakt voor de eerste dagen
van dat plantje --- niet voor ons, al leiden wij het om.

\textbf{2.} In de eigen groene \omterm{def:g1:plants-around-us:parts}{bladeren} van de tarweplanten, uit water en
\omterm{prop:g4:what-plants-need:needs}{minerale zouten} van de bodem en koolstofdioxide uit de lucht, met licht
als \omterm{prop:g3:balanced-diet:jobs}{energie}.

\textbf{3.} Gezaaid doorloopt een korrel haar \omterm{def:g2:born-grow-die:cycle}{levenscyclus}: \omterm{def:g2:from-seed-to-plant:germination}{ontkieming},
\omterm{def:g1:plants-around-us:plant}{plant}, nieuw graan. Verkocht en gemalen eindigt haar \omterm{def:g2:born-grow-die:cycle}{levenscyclus} en
wordt haar \omterm{def:g3:animals-through-seasons:active}{voedselvoorraad} de onze --- lunch ingepakt voor een plantje,
opgegeten door een mens.

\textbf{4.} Geef terug wat je leent: de verterende mest vult de \omterm{prop:g4:what-plants-need:needs}{minerale
zouten} aan die de oogst heeft afgevoerd --- de behoefte aan \omterm{prop:g4:what-plants-need:needs}{minerale
zouten} wordt zo gediend.

\textbf{5.} Omdat \omterm{ex:g6:producing-our-food:bread}{gist} een levende werker is: warm eet en werkt ze snel;
koud ligt ze stil --- de regel van de compostploeg, op bakkerijschaal.

\textbf{6.} Ze voedde zich met de suikers van de \omterm{def:g3:flowers-fruits-seeds:parts}{bloem} en gaf gas af; het
rekbare deeg ving het gas als bellen, en de bellen verdubbelden het deeg.

\textbf{7.} Het koude raam legde de \omterm{ex:g6:producing-our-food:bread}{gist} stil --- nauwelijks gegeten,
nauwelijks gas. De twee groepjes deden per ongeluk een eerlijke proef met
warmte: hetzelfde deeg, dezelfde dag, één omstandigheid anders, twee
rijzingen om te vergelijken.

\textbf{8.} De hitte heeft de \omterm{ex:g6:producing-our-food:bread}{gist} gedood en zo de dienst beëindigd; het
gestolde kruim houdt de gaten van de bellen vast --- de handtekening van
de beroepsbevolking, ter plekke gebakken.

\textbf{9.} Snee uit brood, brood uit deeg, deeg uit \omterm{def:g3:flowers-fruits-seeds:parts}{bloem}, \omterm{def:g3:flowers-fruits-seeds:parts}{bloem} uit
graan, graan uit de \omterm{def:g1:plants-around-us:parts}{bladeren} van de tarweplant --- daar gefabriceerd uit
minerale grondstoffen met licht: het spoor eindigt in de zon van de
juliakker.

\textbf{10.} Snee: één verdieping (van \omterm{def:g5:food-webs:roles}{producent} naar tafel, met een
microbenafwerking). Boter: twee verdiepingen (gras naar koe naar room).
Ham: twee verdiepingen met de grootste tol --- het varken verbrandde en
liet vallen het meeste van de materie van zijn voer voordat er iets ham
werd. De snee is het goedkoopst in materie; de ham het duurst.

\textbf{11.} De \omterm{def:g6:decomposers-and-soil:decomposers}{afbrekers} van de compost --- wormen, pissebedden,
\omterm{def:g6:decomposers-and-soil:decomposers}{schimmels}, onzichtbare ploegen --- halen de \omterm{def:g6:plants-are-producers:matter}{organische materie} van het
deeg uit elkaar tot \omterm{def:g6:decomposers-and-soil:soil}{humus} en \omterm{prop:g4:what-plants-need:needs}{minerale zouten}: één slag verder van voedsel
naar de bodemvruchtbaarheid van volgend jaar.

\textbf{12.} Bijvoorbeeld: ``Ons brood is door een \omterm{def:g5:food-webs:roles}{producent} uit \omterm{def:g6:plants-are-producers:matter}{minerale
materie} en licht gebouwd, door een microbenbevolking opgetild, en had
helemaal geen consumentenverdieping nodig --- al hadden de broodjes ham
ernaast er wel een nodig.''
\end{solution}
